\documentclass[11pt,a4paper]{article}
\usepackage[utf8]{inputenc}
\usepackage[T1]{fontenc}
\usepackage[spanish]{babel}
\usepackage{amsmath,amssymb,mathtools}
\usepackage{geometry}
\usepackage{booktabs}
\usepackage{enumitem}
\usepackage{hyperref}
\usepackage{xcolor}
\usepackage{listings}
\geometry{margin=2.4cm}
\hypersetup{colorlinks=true,linkcolor=blue,urlcolor=blue,citecolor=blue}
\lstset{basicstyle=\ttfamily\small,frame=single,breaklines=true}

\title{Nota de trabajo: funcionamiento matemático y algorítmico de Stockfish}
\author{}
\date{Septiembre de 2026}

\begin{document}
\maketitle

\begin{abstract}
Stockfish es un motor de ajedrez basado en búsqueda selectiva sobre un árbol de juego, con esquema negamax/alpha--beta y técnicas de Principal Variation Search, profundización iterativa, ventanas de aspiración, tabla de transposición, ordenamiento de jugadas, reducciones, extensiones y podas. La evaluación estática moderna se realiza mediante NNUE (Efficiently Updatable Neural Network). Esta nota presenta una formulación matemática compacta del sistema, distinguiendo entre el algoritmo ideal de minimax y las aproximaciones selectivas que permiten alcanzar profundidades prácticas muy grandes.
\end{abstract}

\section{Stockfish no es un único algoritmo}

Una descripción correcta de Stockfish debe separar al menos cuatro subsistemas:
\begin{enumerate}[label=\arabic*.]
    \item representación eficiente de posiciones y generación de jugadas;
    \item búsqueda adversarial de tipo negamax/alpha--beta;
    \item evaluación estática NNUE;
    \item infraestructura de aceleración: tabla de transposición, heurísticas de ordenamiento, poda selectiva, búsqueda paralela y gestión de tiempo.
\end{enumerate}

En términos conceptuales, Stockfish aproxima la solución de
\[
    m^* = \arg\max_{m\in M(s)} V(T(s,m)),
\]
donde $s$ es la posición actual, $M(s)$ el conjunto de jugadas legales, $T(s,m)$ la posición resultante y $V$ el valor minimax de esa posición. El problema es que $V$ no puede calcularse exhaustivamente a profundidad útil, de modo que se reemplaza por una búsqueda selectiva truncada más una evaluación estática.

\section{Minimax y formulación negamax}

Para un juego de suma cero, sea $V_d(s)$ el valor de una posición con horizonte de búsqueda $d$. En forma minimax,
\[
V_d(s)=
\begin{cases}
E(s), & d=0,\\
\max\limits_{m\in M(s)} -V_{d-1}(T(s,m)), & d>0,
\end{cases}
\]
donde $E(s)$ es una evaluación desde el punto de vista del jugador al turno.

Esta forma se denomina \emph{negamax}. Aprovecha la identidad de suma cero
\[
V(s,\text{blancas})=-V(s,\text{negras}).
\]
Por ello no es necesario implementar por separado nodos MAX y MIN.

\section{Poda alpha--beta}

El algoritmo alpha--beta mantiene una ventana $[\alpha,\beta)$ de valores relevantes. Una forma idealizada es:

\begin{lstlisting}
search(s, d, alpha, beta):
    if d <= 0:
        return qsearch(s, alpha, beta)

    best = -infinity
    for m in ordered_moves(s):
        v = -search(T(s,m), d-1, -beta, -alpha)
        best = max(best, v)
        alpha = max(alpha, v)
        if alpha >= beta:
            break
    return best
\end{lstlisting}

Si una rama demuestra que
\[
    v \ge \beta,
\]
se produce un \emph{beta cutoff}: las restantes jugadas no pueden modificar la decisión del nodo padre y se omiten.

Con ordenamiento perfecto, alpha--beta reduce idealmente el coste aproximado desde
\[
O(b^d)
\]
a
\[
O(b^{d/2}),
\]
donde $b$ es el factor medio de ramificación y $d$ la profundidad. En la práctica el beneficio depende fuertemente del orden en que se examinan las jugadas.

\section{Principal Variation Search}

Stockfish distingue la línea principal (PV) de las ramas laterales. Una vez encontrada una jugada razonablemente buena, muchas jugadas siguientes se prueban con una ventana casi nula,
\[
[\alpha,\alpha+\varepsilon),
\]
en vez de una ventana completa. En aritmética discreta puede pensarse como una búsqueda de ventana nula
\[
[\alpha,\alpha+1).
\]

Si la búsqueda estrecha demuestra que una jugada no supera a $\alpha$, se descarta barato. Si supera el límite, se repite con una ventana más amplia. Esta idea constituye el núcleo de Principal Variation Search (PVS).

\section{Profundización iterativa}

Stockfish no comienza directamente a profundidad $D$. Ejecuta sucesivamente
\[
1,2,3,\ldots,D.
\]
Esto parece redundante, pero las búsquedas anteriores producen información crítica:
\begin{itemize}
    \item mejor jugada probable;
    \item línea principal;
    \item valores para la tabla de transposición;
    \item estadísticas históricas para ordenar jugadas;
    \item estimación de estabilidad de la decisión.
\end{itemize}

El coste de las iteraciones poco profundas es pequeño frente al coste de la última, mientras que el mejor ordenamiento reduce fuertemente el árbol final.

\section{Ventanas de aspiración}

Si en la iteración anterior el mejor valor fue $v_{d-1}$, Stockfish puede iniciar la nueva búsqueda con
\[
    \alpha = v_{d-1}-\delta, \qquad
    \beta = v_{d-1}+\delta,
\]
para un $\delta$ pequeño.

Si el resultado cae dentro de la ventana, la búsqueda es más eficiente. Si ocurre
\[
V_d(s)\le \alpha
\]
(fail-low) o
\[
V_d(s)\ge \beta
\]
(fail-high), la ventana se ensancha y se repite la búsqueda.

\section{Tabla de transposición}

Muchas secuencias distintas conducen a la misma posición. Stockfish asocia a cada posición una clave hash de Zobrist $K(s)$ y almacena aproximadamente
\[
\operatorname{TT}[K(s)] = (d,v,B,m,\ldots),
\]
donde $d$ es la profundidad, $v$ el valor, $B$ el tipo de cota y $m$ una jugada recomendada.

Los tipos conceptuales de cota son:
\begin{align*}
B &= \text{EXACT}, && V(s)=v,\\
B &= \text{LOWER}, && V(s)\ge v,\\
B &= \text{UPPER}, && V(s)\le v.
\end{align*}

Si la profundidad almacenada y el tipo de cota son suficientes para la ventana actual, puede evitarse una nueva exploración. Incluso cuando no permite cortar, la jugada almacenada suele probarse primero, mejorando el ordenamiento.

\section{Ordenamiento de jugadas}

La eficiencia de alpha--beta depende más del ordenamiento que del número bruto de jugadas. Stockfish combina, entre otras fuentes de información:
\begin{itemize}
    \item jugada de la tabla de transposición;
    \item capturas prometedoras;
    \item información SEE (Static Exchange Evaluation);
    \item estadísticas históricas;
    \item continuation history;
    \item información aprendida durante la propia búsqueda.
\end{itemize}

En forma abstracta, cada movimiento recibe un puntaje de prioridad
\[
P(m\mid s)=w_1 H(m,s)+w_2 C(m,s)+w_3 T(m,s)+\cdots,
\]
y se exploran antes las jugadas con mayor $P$.

No debe interpretarse $P$ como evaluación ajedrecística del movimiento. Es, ante todo, una estimación de qué jugada tiene mayor probabilidad de producir pronto un cutoff.

\section{Late Move Reductions}

Una de las técnicas selectivas más importantes es Late Move Reduction (LMR). Si una jugada aparece tarde en el orden, es tranquila y el contexto no sugiere una táctica crítica, inicialmente se busca a profundidad reducida:
\[
    d' = d-1-r(d,n,\text{contexto}),
\]
donde $n$ es aproximadamente el número ordinal de la jugada y $r$ es una reducción creciente con $d$ y $n$.

Una forma conceptual típica es
\[
    r(d,n) \approx a+b\ln d\,\ln n,
\]
aunque la implementación real utiliza funciones y correcciones empíricamente ajustadas.

Si la búsqueda reducida devuelve un valor inesperadamente bueno,
\[
    v > \alpha,
\]
la jugada puede reexaminarse con menor reducción o profundidad completa.

La lógica probabilística implícita es: las jugadas tardías, después de un buen ordenamiento, tienen baja probabilidad de ser la mejor jugada. Se invierte menos cómputo primero y se profundiza sólo cuando la evidencia contradice esa hipótesis.

\section{Podas selectivas}

Stockfish incorpora varias podas y búsquedas de verificación. Su principio general es evitar invertir profundidad completa en ramas con baja probabilidad de alterar la decisión.

Una representación abstracta es:
\[
\text{buscar rama a profundidad completa}
\quad\Longleftrightarrow\quad
\Pr(\text{rama cambia la decisión}\mid\mathcal I) > \tau,
\]
donde $\mathcal I$ representa información disponible y $\tau$ un umbral implícito.

Entre las familias usadas históricamente y en versiones modernas del motor aparecen:
\begin{itemize}
    \item poda por evaluación estática;
    \item búsquedas de ventana nula;
    \item null-move pruning y verificaciones relacionadas;
    \item futility pruning;
    \item razoring o mecanismos equivalentes según versión;
    \item ProbCut o técnicas equivalentes según versión;
    \item reducciones dependientes de historial y contexto.
\end{itemize}

Los detalles exactos cambian con frecuencia porque Stockfish evoluciona mediante experimentación A/B distribuida en Fishtest. Por ello conviene separar el principio matemático estable de los coeficientes y reglas concretas de una revisión determinada.

\section{Quiescence search}

Detener la búsqueda exactamente en $d=0$ puede producir el \emph{efecto horizonte}. Por ejemplo, una pieza puede parecer ganada justo antes de una recaptura obvia.

En lugar de aplicar inmediatamente $E(s)$, Stockfish entra en una búsqueda de quietud:
\[
Q(s,\alpha,\beta).
\]
Se calcula primero un valor estático $e=E(s)$ y luego se exploran selectivamente jugadas tácticas relevantes hasta alcanzar posiciones suficientemente estables.

Conceptualmente:
\[
Q(s)=\max\left(E(s),\max_{m\in T(s)} -Q(T(s,m))\right),
\]
donde $T(s)$ es un subconjunto de movimientos tácticos apropiados.

\section{Evaluación NNUE}

Las versiones modernas de Stockfish utilizan NNUE en lugar de la antigua evaluación artesanal clásica. NNUE significa \emph{Efficiently Updatable Neural Network}.

Sea $x(s)$ un vector disperso de características de la posición. Una red feed-forward simplificada puede expresarse como
\begin{align*}
h_1 &= \phi(W_1x+b_1),\\
h_2 &= \phi(W_2h_1+b_2),\\
y   &= W_3h_2+b_3,
\end{align*}
donde $\phi$ es una activación saturada del tipo ClippedReLU.

Una forma idealizada es
\[
\phi(z)=\min(1,\max(0,z)),
\]
aunque la implementación usa aritmética cuantizada entera y escalas específicas.

La arquitectura actual de Stockfish incluye una transformación de características amplia seguida de capas densas pequeñas. En el código oficial de septiembre de 2026 aparecen dimensiones principales de transformación $1024$ y capas densas de $32$ y $32$ unidades, con activaciones ClippedReLU y SqrClippedReLU.

\subsection{Actualización incremental}

La ventaja esencial de NNUE no es solamente la red sino que la primera transformación puede actualizarse incrementalmente. Si entre $s$ y $s'$ cambian unas pocas características,
\[
 x(s')=x(s)+\Delta x,
\]
entonces el acumulador lineal satisface
\[
 a(s')=W x(s')+b = a(s)+W\Delta x.
\]

Como $\Delta x$ es muy disperso, no se recalcula la gran matriz completa. Se suman o restan solamente las columnas correspondientes a características activadas o desactivadas.

Esto hace posible evaluar millones de nodos con una red neuronal usando CPU convencional.

\section{Escalado de la evaluación}

La salida NNUE no se utiliza necesariamente de manera totalmente aislada. La implementación actual aplica ajustes dependientes del estado de búsqueda y de la posición. Por ejemplo, el código oficial de septiembre de 2026 combina la evaluación NNUE con una medida material simple para estimar alineación y aplica además escalado por material y por el contador de la regla de cincuenta movimientos.

Conceptualmente puede escribirse
\[
E_{\mathrm{final}}(s)
= S\left(E_{\mathrm{NNUE}}(s),M(s),R_{50}(s),O(s)\right),
\]
donde $M$ representa información material, $R_{50}$ el contador de la regla de 50 jugadas y $O$ un término dependiente del contexto de búsqueda.

Debe evitarse interpretar la salida de Stockfish como una probabilidad directa. El valor interno sirve principalmente para ordenar y comparar posiciones dentro de la búsqueda.

\section{Mate y valores terminales}

Las posiciones terminales tienen prioridad sobre NNUE. Si el jugador al turno está en mate a distancia $k$, se utilizan valores especiales de mate cuya magnitud domina cualquier evaluación posicional.

Una representación típica es
\[
V_{mate}(k)=M-k,
\]
para una victoria, y
\[
V_{mated}(k)=-M+k,
\]
para una derrota, con $M$ suficientemente grande.

Esto hace que entre dos mates Stockfish prefiera el mate más rápido, y cuando está perdido tienda a retrasarlo.

\section{Syzygy tablebases}

Cuando una posición entra en el dominio cubierto por las tablebases configuradas, no es necesario aproximar el resultado: puede consultarse información exacta de finales. En posiciones externas a la tablebase, las consultas durante la búsqueda pueden permitir cortes o selección de jugadas que preservan victoria o tablas.

En ese caso el motor está combinando tres niveles de conocimiento:
\[
\text{búsqueda} + \text{NNUE} + \text{información exacta de tablebase}.
\]

\section{Paralelismo}

Stockfish utiliza múltiples hilos que exploran el árbol y comparten información, especialmente la tabla de transposición. La documentación denomina a este enfoque \emph{Lazy SMP}.

No debe imaginarse como una partición perfecta del árbol. Diferentes hilos pueden explorar variantes parcialmente redundantes, pero esa diversidad produce información que mejora los cutoffs de los demás hilos mediante la TT.

Una forma conceptual del beneficio es reducir el tiempo efectivo para obtener una cierta cantidad de información útil:
\[
T_p < T_1,
\]
pero sin esperar escalado lineal ideal
\[
T_p \neq \frac{T_1}{p}
\]
debido a redundancia, sincronización, ancho de banda de memoria y colisiones en estructuras compartidas.

\section{Gestión de tiempo}

Stockfish no fija simplemente una profundidad. En una partida real decide cuánto tiempo invertir en función de, entre otros elementos:
\begin{itemize}
    \item tiempo restante;
    \item incremento;
    \item estabilidad de la mejor jugada;
    \item variación del valor entre iteraciones;
    \item esfuerzo relativo invertido en la mejor línea;
    \item límites UCI solicitados.
\end{itemize}

Una representación conceptual es
\[
 t_{\mathrm{move}}
 = t_{\mathrm{base}}\cdot
 f_{\mathrm{stability}}\cdot
 f_{\mathrm{score}}\cdot
 f_{\mathrm{effort}}.
\]

Si la mejor jugada cambia repetidamente o la evaluación se vuelve inestable, suele justificarse gastar más tiempo. Si una única jugada permanece claramente dominante, puede detenerse antes.

\section{Flujo global simplificado}

El proceso completo puede resumirse como:
\begin{enumerate}
    \item generar las jugadas legales de la raíz;
    \item iniciar profundización iterativa;
    \item construir una ventana de aspiración alrededor del resultado previo;
    \item consultar la tabla de transposición;
    \item obtener evaluación estática NNUE cuando corresponde;
    \item ordenar movimientos usando TT, capturas, SEE e historiales;
    \item buscar la primera jugada con ventana amplia;
    \item buscar muchas jugadas posteriores con PVS y reducciones LMR;
    \item aplicar podas selectivas cuando la evidencia indica que una rama no puede modificar la decisión;
    \item ampliar o repetir la búsqueda si una búsqueda reducida produce un resultado sorprendente;
    \item al llegar al horizonte, entrar en quiescence search;
    \item almacenar resultados y cotas en la TT;
    \item repetir a mayor profundidad hasta que el gestor de tiempo ordene detenerse;
    \item devolver la mejor jugada de la última iteración fiable.
\end{enumerate}

\section{Interpretación matemática}

Puede pensarse Stockfish como un estimador adaptativo del valor minimax:
\[
\widehat V(s;C)=\mathcal S\bigl(s,C,E_{\mathrm{NNUE}},\mathcal H,\operatorname{TT}\bigr),
\]
donde $C$ es el presupuesto computacional, $\mathcal S$ la búsqueda selectiva y $\mathcal H$ el conjunto de heurísticas de ordenamiento y reducción.

Cuando $C$ aumenta, idealmente
\[
\widehat V(s;C) \longrightarrow V(s),
\]
pero en ajedrez práctico no se alcanza el árbol completo. La fuerza proviene de asignar el presupuesto de forma extremadamente no uniforme: mucha búsqueda en ramas plausibles y muy poca en ramas estadísticamente improbables.

Por eso la profundidad UCI informada por Stockfish no equivale a ``todas las variantes fueron examinadas hasta $d$''. El árbol es fuertemente selectivo y algunas líneas alcanzan profundidades considerablemente mayores que otras.

\section{Diferencia conceptual con un enfoque de refuerzo tipo Bellman}

Una ecuación de Bellman para un agente aprendido podría escribirse
\[
Q(s,a) \leftarrow Q(s,a)+\eta\left[r+\gamma\max_{a'}Q(s',a')-Q(s,a)\right].
\]

Stockfish no basa su decisión principal en una tabla $Q(s,a)$ aprendida de esa manera. Su mecanismo central es una búsqueda explícita del árbol en tiempo de decisión. NNUE aporta una estimación $E(s)$ en las hojas o nodos estáticos, pero el valor final de una jugada surge de propagar resultados de búsqueda adversarial.

En forma esquemática:
\[
\text{Stockfish:}\quad a^*\approx \arg\max_a\left[-\operatorname{Search}(T(s,a))\right],
\]
mientras que un agente puramente basado en $Q$ usaría
\[
\text{Q-learning:}\quad a^*=\arg\max_a Q(s,a).
\]

Esta distinción es fundamental al integrar ideas de Stockfish en otros motores: copiar una evaluación estática no reproduce Stockfish; gran parte de su fuerza reside en cómo distribuye selectivamente la búsqueda.

\section{Conclusiones}

La descripción más compacta y útil es la siguiente:
\begin{quote}
Stockfish es un buscador negamax alpha--beta altamente selectivo, organizado mediante profundización iterativa y PVS, que reutiliza resultados en una tabla de transposición, concentra el esfuerzo mediante ordenamiento, podas y reducciones, estabiliza las hojas tácticas con quiescence search y utiliza NNUE como función de evaluación estática incremental. La decisión final es el resultado de la interacción entre evaluación y búsqueda, no de ninguna de ellas por separado.
\end{quote}

Desde el punto de vista matemático, la principal ingeniería no consiste en cambiar la ecuación minimax, sino en estimar qué partes del árbol merecen cómputo completo y cuáles pueden ser aproximadas, reducidas o descartadas sin perder fuerza estadística.

\section*{Referencias técnicas consultadas}
\begin{itemize}
    \item Stockfish Developers, \emph{Stockfish source code}, especialmente \texttt{src/search.cpp}, \texttt{src/evaluate.cpp} y \texttt{src/nnue/nnue\_architecture.h}:\\
    \url{https://github.com/official-stockfish/Stockfish}
    \item Stockfish Documentation, \emph{Advanced Topics}:\\
    \url{https://official-stockfish.github.io/docs/stockfish-wiki/Advanced-topics.html}
    \item Stockfish Documentation, \emph{Terminology}:\\
    \url{https://official-stockfish.github.io/docs/stockfish-wiki/Terminology.html}
    \item Stockfish Developers, \emph{NNUE PyTorch documentation}:\\
    \url{https://github.com/official-stockfish/nnue-pytorch}
\end{itemize}

\end{document}
